\begin{graphicspathcontext}{{./chapters/ai-health/imgs/},{./chapters/ai-health/imgs/auto/},{./chapters/ai/imgs/},{./chapters/ai/imgs/auto/},\old}

% -----------------------------------------------------------------------
\begin{frame}[t]{{L'IA en santé :} Une histoire riche}
	\begin{stabularx}{X|X}
	\tabularhead{Période}{Avancées majeures} \\
	Années 70 & \textbf{MYCIN} : premier système expert de diagnostic médical (Stanford) \\
	\hline
	Années 90 & Systèmes d'aide au diagnostic par imagerie (CAD) \\
	\hline
	2012--2016 & \textit{Deep learning} : détection de la rétinopathie diabétique (Google DeepMind) \\
	\hline
	2017--2020 & Premières approbations FDA pour des IA de diagnostic radiologique \\
	\hline
	2022 & \textbf{AlphaFold 2} : prédiction de la structure de toutes les protéines connues \\
	\hline
	2023--2025 & LLMs et agents conversationnels en support clinique \\
	\end{stabularx}
\end{frame}

% -----------------------------------------------------------------------
\figureslide{{L'IA en santé :} Vue d'ensemble des applications}{ai_health_applications}

% -----------------------------------------------------------------------
\begin{frame}[t]{{L'IA en santé :} De nouvelles perspectives}
	\alertbox{L'essor de l'IA générative et des LLMs ouvre aujourd'hui de nouvelles voies}
	\vspace{.25cm}
	\begin{rightarrowsequence}
	\arrow{Assistance au diagnostic}
	\arrow[bg=CIADlightgray,fg=black]{Analyse d'images médicales, aide à la décision clinique, interprétation d'examens}
	\end{rightarrowsequence}
	\vspace{.2cm}
	\begin{rightarrowsequence}
	\arrow{Support thérapeutique}
	\arrow[bg=CIADlightgray,fg=black]{Agents conversationnels pour le suivi des patients et l'observance}
	\end{rightarrowsequence}
	\vspace{.2cm}
	\begin{rightarrowsequence}
	\arrow{Optimisation organisationnelle}
	\arrow[bg=CIADlightgray,fg=black]{Planification, gestion des flux patients, logistique médicale}
	\end{rightarrowsequence}
	\vspace{.2cm}
	\begin{rightarrowsequence}
	\arrow{Recherche biomédicale}
	\arrow[bg=CIADlightgray,fg=black]{Découverte de médicaments, médecine de précision, génomique}
	\end{rightarrowsequence}
\end{frame}

% -----------------------------------------------------------------------
\begin{frame}[t]{{L'IA en santé :} Assistance au diagnostic}
	\footnotesize
	\alertbox*{Plus de 500 outils d'IA approuvés par la FDA pour la radiologie (2024)}
	\begin{columns}
		\begin{column}{.48\linewidth}
			\begin{rightanchorblock}{Imagerie médicale}{1}
			Radiologie, IRM, scanner : détection de tumeurs, fractures, pathologies pulmonaires
			\end{rightanchorblock}
			\begin{rightanchorblock}{Anatomopathologie}{2}
			Analyse de lames numérisées pour la détection de cellules cancéreuses
			\end{rightanchorblock}
			\begin{rightanchorblock}{Cardiologie}{3}
			Interprétation automatique des ECG, détection de fibrillations atriales
			\end{rightanchorblock}
		\end{column}
		\begin{column}{.52\linewidth}
			\vspace{.1cm}
			\begin{rightarrowsequence}
			\arrow{Dermatologie}
			\arrow[bg=CIADlightgray,fg=black]{Détection de mélanomes à partir de photos cutanées (Google DermAssist)}
			\end{rightarrowsequence}
			\vspace{.2cm}
			\begin{rightarrowsequence}
			\arrow{Ophtalmologie}
			\arrow[bg=CIADlightgray,fg=black]{Dépistage de la rétinopathie diabétique — performances comparables à l'expert}
			\end{rightarrowsequence}
		\end{column}
	\end{columns}
	\normalsize
\end{frame}

% -----------------------------------------------------------------------
\begin{frame}[t]{{L'IA en santé :} Utilisation thérapeutique des LLMs}
	\footnotesize
	\begin{columns}
		\begin{column}{.55\linewidth}
			\begin{rightanchorblock}{Arrêt du tabac}{1}
			Étude (2024) : les agents conversationnels IA améliorent les taux d'abstinence à 3 mois vs. interventions classiques, avec une disponibilité 24h/24
			\end{rightanchorblock}
			\begin{rightanchorblock}{Santé mentale}{2}
			\textit{Woebot}, \textit{Wysa} : suivi de l'anxiété et de la dépression, TCC numérique, détection de signaux de crise
			\end{rightanchorblock}
			\begin{rightanchorblock}{Observance thérapeutique}{3}
			Rappels personnalisés, explication des traitements, suivi des effets indésirables, éducation thérapeutique
			\end{rightanchorblock}
		\end{column}
		\begin{column}{.45\linewidth}
			\alertbox{Les agents conversationnels \Emph{complètent} la prise en charge --- ils ne remplacent pas le médecin}
			\vspace{.3cm}
			\alertbox*{Enjeu clé : validation clinique et encadrement réglementaire (dispositif médical logiciel)}
		\end{column}
	\end{columns}
	\normalsize
\end{frame}

% -----------------------------------------------------------------------
\begin{frame}[t]{{L'IA en santé :} Optimisation et planification}
	\begin{stabularx}{X|X|X}
	\tabularhead{Domaine}{Problème}{Apport de l'IA} \\
	Planning du personnel & Plannings infirmiers multi-contraintes & Optimisation combinatoire, saisie en langage naturel \\
	\hline
	Blocs opératoires & Ordonnancement des interventions & Minimisation des annulations, prédiction des durées \\
	\hline
	Flux patients & Anticipation des admissions aux urgences & Modèles prédictifs pour les pics d'activité \\
	\hline
	Logistique médicale & Approvisionnement en médicaments & Prévision des stocks, réduction des ruptures \\
	\end{stabularx}
\end{frame}


% -----------------------------------------------------------------------
\figureslide{{L'IA en santé :} Des défis à résoudre}{ai_health_challenges}

% -----------------------------------------------------------------------
\begin{frame}[t]{{L'IA en santé :} Données de santé et confidentialité}
	\footnotesize
	\begin{columns}
		\begin{column}{.5\linewidth}
			\alertbox*{Les données de santé sont parmi les plus sensibles}
			\vspace{.2cm}
			\begin{rightanchorblock}{RGPD (art. 9)}{1}
				Données de santé = catégorie particulière. Traitement soumis à des conditions strictes et au consentement explicite
			\end{rightanchorblock}
			\begin{rightanchorblock}{Certification HDS}{2}
				Hébergement des Données de Santé : certification obligatoire en France pour tout hébergement de données patients
			\end{rightanchorblock}
		\end{column}
		\begin{column}{.5\linewidth}
			\begin{rightanchorblock}{Souveraineté numérique}{3}
				Où les données sont-elles hébergées ? Le CLOUD Act américain permet l'accès aux données même hébergées en Europe
			\end{rightanchorblock}
			\begin{rightanchorblock}{Limites de l'anonymisation}{4}
				La ré-identification reste possible même sur des données « anonymisées » --- problème de désanonymisation
			\end{rightanchorblock}
		\end{column}
	\end{columns}
	\normalsize
\end{frame}

% -----------------------------------------------------------------------
\begin{frame}[t]{{L'IA en santé :} L'EU AI Act et le secteur médical}
	\begin{columns}
		\begin{column}{.45\linewidth}
			\alertbox{La majorité des IA médicales sont classées \Emph{Risque Élevé} par l'EU AI Act}
			\vspace{.3cm}
			\begin{rightanchorblock}{Exigences}{1}
			Évaluation de conformité, journalisation, supervision humaine, transparence sur les données d'entraînement
			\end{rightanchorblock}
		\end{column}
		\begin{column}{.55\linewidth}
			\begin{stabularx}{X|X}
			\tabularhead{Catégorie IA}{Exemples en santé} \\
			\textbf{Interdit} & Évaluation du risque de récidive psychiatrique \\
			\hline
			\textbf{Risque Élevé} & Diagnostic assisté, priorisation aux urgences \\
			\hline
			\textbf{Risque Limité} & Chatbots de santé généralistes, information patient \\
			\hline
			\textbf{Risque Minimal} & Rappels médicaments, calcul d'IMC \\
			\end{stabularx}
		\end{column}
	\end{columns}
\end{frame}

\end{graphicspathcontext}

\endinput
